\chapter{Flushing}

\section*{Definition}
Sudden, transient erythema of the skin, most commonly on the face, neck, and upper chest, caused by widening of blood vessels of cutaneous blood vessels; may be episodic or persistent depending on causes.

\section*{What Happens in Your Body}
Flushing results from relaxation of cutaneous vascular smooth muscle mediated by vasoactive substances including histamine, serotonin, prostaglandins, nitric oxide, and catecholamines. The blush area corresponds to the distribution of the trigeminal nerve (V1-V2) and cervical dermatomes where thermoregulatory widening of blood vessels is most active.

\section*{Causes}
\begin{itemize}
  \item Physiologic: Menopausal hot flushes, emotional blushing, heat exposure, exercise, spicy foods (capsaicin), alcohol (especially in Asian populations with ALDH2 deficiency)
  \item Medications: Niacin (chemical messengers that cause inflammation and pain-mediated), calcium channel blockers, nitrates, tamoxifen, sildenafil, steroid withdrawal, vancomycin (red man syndrome)
  \item Neoplastic: Carcinoid syndrome (flushing with diarrhea, wheezing, right-sided heart failure), pheochromocytoma (paroxysmal hypertension, palpitations), medullary thyroid carcinoma, VIPoma
  \item Endocrine: Hyperthyroidism, systemic mastocytosis, rosacea (persistent facial erythema with telangiectasias)
\end{itemize}

\section*{When to See a Doctor}
See your doctor if:
\begin{itemize}
  \item Flushing with diarrhea, wheezing, or heart murmur suggests carcinoid syndrome
  \item Paroxysmal flushing with severe hypertension, headache, and palpitations
  \item Persistent or progressive flushing unresponsive to trigger avoidance
  \item Flushing accompanied by unexplained weight loss, fevers, or night sweats
\end{itemize}

\section*{Related Symptoms}
\begin{itemize}
  \item Palpitations
  \item Diarrhea
  \item Headache
  \item Diaphoresis
  \item Telangiectasias
\end{itemize}
\textbf{Important:} Niacin-induced flushing is common but benign; use of aspirin 30 minutes before niacin can reduce the vasodilatory chemical messengers that cause inflammation and pain response.

\section*{Self-Care / Home Management}
\begin{itemize}
  \item Identify and avoid triggers (spicy foods, alcohol, hot beverages, stressful situations).
  \item Use cool compresses or a fan during flushing episodes for symptomatic relief.
  \item For menopausal hot flushes, dress in layers, maintain a cool environment, and consider soy isoflavones or black cohosh (variable evidence).
  \item Avoid self-medicating with prescription medications without medical guidance.
\end{itemize}

\section*{How Your Doctor Will Evaluate You}
\begin{itemize}
  \item 24-hour urinary 5-HIAA (5-hydroxyindoleacetic acid) for carcinoid syndrome workup.
  \item Plasma metanephrines or 24-hour urinary catecholamines for pheochromocytoma.
  \item Serum tryptase, urinary histamine metabolites for mastocytosis.
  \item Thyroid function tests, menopause confirmation (FSH, estradiol) as clinically indicated.
\end{itemize}

\section*{Treatment Options}
\begin{itemize}
  \item Somatostatin analogs (octreotide) for symptomatic carcinoid syndrome.
  \item Alpha-blockade (phenoxybenzamine) followed by beta-blockade for pheochromocytoma; surgical resection.
  \item Hormone replacement therapy or SSRIs/SNRIs for menopausal vasomotor symptoms.
  \item Avoidance of triggers; niacinamide (no-flush niacin) alternative for dyslipidemia treatment.
\end{itemize}


\section*{References}
\begin{thebibliography}{99}
\bibitem{izikson2006}
Izikson L, English JC III, Zirwas MJ. The flushing patient: differential diagnosis, workup, and treatment. \textit{J Am Acad Dermatol}. 2006;55(2):193--208.
\bibitem{mohammed2011}
Mohammed K, Ghosh S. Management of flushing and hot flushes. \textit{Br J Hosp Med}. 2011;72(10):M156--M159.
\bibitem{altschuler2002}
Altschuler EL, Kast RE. The flushing patient: a review. \textit{Int J Dermatol}. 2002;41(4):193--200.
\bibitem{ray1996}
Ray WA. Flushing and cardiovascular disease. \textit{JAMA}. 1996;276(5):375--376.
\bibitem{lonardi2011}
Lonardi S, Bortolami A, Stefani M, et al. Carcinoid syndrome and flushing. \textit{Ann Oncol}. 2011;22(suppl 5):v74--v78.
\bibitem{wilkin1981}
Wilkin JK. Flushing reactions: consequences and mechanisms. \textit{Ann Intern Med}. 1981;95(4):468--476.
\end{thebibliography}
\clearpage
